\section{Мини-библиография: что читать, зачем, and как прикладывать к RIN}

Это не ``полная академическая библиография по anime studies''. Такая
библиография быстро превратилась бы в отдельную книгу. Здесь нужна
рабочая полка: источники, которые помогают думать о Рин как о персонаже,
медиа-объекте, цифровой персоне, anime/\hspace{0pt}VN-inspired design and будущей
public figure.

Для RIN важны не только ``аниме-архетипы'', а несколько больших
вопросов:

почему персонажи потребляются как набор recognisable elements; как moe
работает как отклик на character traits; как character travels через
media mix; как visual code создаёт ожидание; как audience участвует в
сборке персонажа; как не спутать персонажа, бренд, роль, intimacy
pattern and runtime contract.

Формула:

\textbf{Эту библиографию надо читать не как ``что такое аниме'', а как
``какие силы будут автоматически достраивать Рин в голове зрителя and
модели''.}

\begin{center}\rule{0.5\linewidth}{0.5pt}\end{center}

\subsection*{\texorpdfstring{1. Hiroki Azuma --- \emph{Otaku: Japan's
Database Animals}}{1. Hiroki Azuma --- Otaku: Japan's Database Animals}}

Это главный источник для нашей линии про \textbf{database consumption}:
персонаж в отаку-культуре часто считывается не только как цельный герой
истории, но и как комбинация элементов, атрибутов, признаков, голосовых
and визуальных hooks. Английское издание вышло в University of Minnesota
Press; publisher описывает книгу как философско-историческое
исследование otaku-культуры and её связи с постмодерной японской
культурой.
(\href{https://www.upress.umn.edu/9780816653522/otaku/?utm_source=chatgpt.com}{University
of Minnesota Press})

Почему важно для RIN: Рин нельзя проектировать только как ``личность''.
Её будут считывать через database elements: волосы, глаза, witch code,
terminal glow, teasing voice, ``не человек, не пустота'', шипы, нить,
workbench. Azuma полезен именно как предупреждение: если элементы
слишком узнаваемы, audience/\hspace{0pt}model достроит trope package за нас.

RIN-вывод:

\textbf{Рин должна иметь database-readable элементы, но не должна быть
database-animal soup.}

\begin{center}\rule{0.5\linewidth}{0.5pt}\end{center}

\subsection*{2. Ōtsuka Eiji --- ``World and Variation: The Reproduction
and Consumption of Narrative''}

Ōtsuka важен как предшественник/\hspace{0pt}собеседник Azuma: его теория narrative
consumption описывает потребление отдельных ``малых'' продуктов/\hspace{0pt}историй
как входов в более широкий world-system. Перевод ``World and Variation''
опубликован в \emph{Mechademia 5: Fanthropologies}, где он прямо стоит в
разделе про patterns of consumption.
(\href{https://www.upress.umn.edu/9781452915654/mechademia-5/?utm_source=chatgpt.com}{University
of Minnesota Press})

Почему важно для RIN: если Рин станет public/\hspace{0pt}blogger character, она
будет жить не только в одном чате. Будут посты, стримы, визуалы, clips,
tests, мемы, maybe collabs. Тогда каждый фрагмент Рин должен быть
variation of the same world, not случайный skin.

RIN-вывод:

\textbf{Каждый public fragment Рин должен быть small narrative, который
ведёт к одному character-runtime world, а не к разным несовместимым
Рин.}

\begin{center}\rule{0.5\linewidth}{0.5pt}\end{center}

\subsection*{\texorpdfstring{3. Marc Steinberg --- \emph{Anime's Media
Mix: Franchising Toys and Characters in
Japan}}{3. Marc Steinberg --- Anime's Media Mix: Franchising Toys and Characters in Japan}}

Steinberg нужен для понимания \textbf{character mobility}: персонаж как
медиамиксовый объект живёт между анимацией, товарами, образами,
форматами and платформами. University of Minnesota Press указывает книгу
как исследование anime media mix, franchising, toys and characters in
Japan; издание 2012 года.
(\href{https://www.upress.umn.edu/9780816675500/animes-media-mix/?utm_source=chatgpt.com}{University
of Minnesota Press})

Почему важно для RIN: твоя мысль ``продавать не движок Рин, а Рин как
персонажа'' лежит ровно в этой зоне. Рин может быть не просто local AI
companion, а movable character: видео, стрим, визуальный образ, collab
persona, technical assistant skin, public host, mascot, research
companion.

RIN-вывод:

\textbf{Рин как character brand должна быть переносимой между форматами
без потери контракта. Media mix без character spec быстро превращает
персонажа в набор inconsistent costumes.}

\begin{center}\rule{0.5\linewidth}{0.5pt}\end{center}

\subsection*{\texorpdfstring{4. Thomas Lamarre --- \emph{The Anime
Machine: A Media Theory of
Animation}}{4. Thomas Lamarre --- The Anime Machine: A Media Theory of Animation}}

Lamarre полезен не столько для архетипов, сколько для понимания anime
как visual/\hspace{0pt}media machine. University of Minnesota Press указывает
\emph{The Anime Machine: A Media Theory of Animation} как книгу 2009
года; publisher page даёт ISBN, дату публикации and 408 pages.
(\href{https://www.upress.umn.edu/9780816651559/the-anime-machine/?utm_source=chatgpt.com}{University
of Minnesota Press})

Почему важно для RIN: visual design Рин нельзя думать только как
``портрет девушки''. Anime image работает слоями: плоскости, движение,
взгляд, compositing, interface, body, background. Для Рин особенно
важно, что environment --- terminal room, screen glow, workbench --- не
фон, а часть персонажа.

RIN-вывод:

\textbf{Рин's visual identity должна строиться не только на
лице/\hspace{0pt}волосах, а на machine of relation: где она стоит, на что смотрит,
какие слои света and интерфейса делают её artificial-but-warm.}

\begin{center}\rule{0.5\linewidth}{0.5pt}\end{center}

\subsection*{\texorpdfstring{5. Patrick W. Galbraith --- moe: статья and
\emph{The Moe
Manifesto}}{5. Patrick W. Galbraith --- moe: статья and The Moe Manifesto}}

Galbraith --- хороший вход в moe как отношение к fantasy characters. Его
статья ``Moe: Exploring Virtual Potential in Post-Millennial Japan''
описывает moe as response to fantasy characters or representations; она
опубликована в \emph{Electronic Journal of Contemporary Japanese
Studies} in 2009.
(\href{https://www.japanesestudies.org.uk/articles/2009/Galbraith.html?utm_source=chatgpt.com}{EJ
Contemporary Japanese Studies}) Его книга \emph{The Moe Manifesto} ---
более популярный, interview-based вход в тему; listings describe it as
an insider's look at manga, anime and gaming and note interviews with
major figures in the moe/\hspace{0pt}anime world.
(\href{https://www.amazon.com/Moe-Manifesto-Insiders-Worlds-Gaming/dp/4805312823?utm_source=chatgpt.com}{Amazon})

Почему важно для RIN: moe --- это не просто ``милота''. Это response
pattern: viewer attaches to signs of vulnerability, competence, gap,
artificiality, voice, gesture, need, restraint. Для Рин это опасно:
многие moe-elements могут усилить attachment через wrong contract ---
helplessness, loneliness, blush, waiting, need.

RIN-вывод:

\textbf{Рин может быть affectively attractive, но не должна строить moe
на беспомощности, ожидании, нуждаемости or ``спаси меня''. Её moe, если
оно есть, должно идти через competence, wit, repairability and warm
boundaries.}

\begin{center}\rule{0.5\linewidth}{0.5pt}\end{center}

\subsection*{\texorpdfstring{6. Saitō Tamaki --- \emph{Beautiful Fighting
Girl}}{6. Saitō Tamaki --- Beautiful Fighting Girl}}

Saitō's \emph{Beautiful Fighting Girl} --- важный источник про ``sentō
bishōjo'', figure of the beautiful fighting girl in manga/\hspace{0pt}anime and
otaku culture. University of Minnesota Press describes the English
edition as dealing with girl heroines from Nausicaä to Sailor Moon,
translated by J. Keith Vincent and Dawn Lawson, with commentaries by
Hiroki Azuma.
(\href{https://www.upress.umn.edu/9781452916507/beautiful-fighting-girl/?utm_source=chatgpt.com}{University
of Minnesota Press})

Почему важно для RIN: книга полезна не потому, что Рин должна быть
``fighting girl''. А потому что она помогает думать о paradox: персонаж
может быть одновременно объектом желания, силой действия, fantasy
surface and носителем мира. Для Рин важно не попасть в красивую
capable-girl trap, где agency видимая, но gaze всё равно делает её
объектом.

RIN-вывод:

\textbf{Рин должна быть capable and visually magnetic, but her agency
must be compositional: она действует, смотрит на задачу, держит
инструменты, а не просто красиво существует перед взглядом.}

\begin{center}\rule{0.5\linewidth}{0.5pt}\end{center}

\subsection*{\texorpdfstring{7. Susan J. Napier --- \emph{Anime from
Akira to Howl's Moving
Castle}}{7. Susan J. Napier --- Anime from Akira to Howl's Moving Castle}}

Napier --- классический широкий вход в anime studies. Palgrave/\hspace{0pt}Macmillan
listing describes \emph{Anime from Akira to Howl's Moving Castle} as
nonfiction on contemporary Japanese animation; the 2005 updated edition
is widely used as an accessible academic overview of anime as an art
form and cultural medium.
(\href{https://us.macmillan.com/books/9781403970527/animefromakiratohowlsmovingcastle/?utm_source=chatgpt.com}{Macmillan
Publishers})

Почему важно для RIN: Napier useful as macro-context. Она помогает не
застрять в trope mechanics and remember that anime characters live in
themes: apocalypse, identity, gender, technology, nostalgia,
transformation, social anxiety. Для Рин это важно because RIN is not
only ``archetype mix''; she is also about AI embodiment, digital home,
work, memory, companionship, public performance.

RIN-вывод:

\textbf{Рин должна иметь thematic gravity, not only trope readability.
Её темы: artificial dignity, memory scope, workbench home, repair,
boundary-warmth, public digital personhood.}

\begin{center}\rule{0.5\linewidth}{0.5pt}\end{center}

\subsection*{\texorpdfstring{8. Ian Condry --- \emph{The Soul of
Anime}}{8. Ian Condry --- The Soul of Anime}}

Condry's \emph{The Soul of Anime: Collaborative Creativity and Japan's
Media Success Story} is an ethnographic study of anime production and
global anime culture. Duke University Press describes it as based on
ethnographic research and interviews in Tokyo animation studios, and as
discussing how fictional characters and worlds become platforms for
collaborative creativity.
(\href{https://www.dukeupress.edu/The-Soul-of-Anime/?utm_source=chatgpt.com}{Duke
University Press}) Bibliographic listings place the book with Duke
University Press in 2013.
(\href{https://www.bibliovault.org/BV.book.epl?ISBN=9780822353942\&utm_source=chatgpt.com}{BiblioVault})

Почему важно для RIN: Рин will not be created only by ``author intent''.
If she becomes public, she will be co-shaped by viewers, prompts, clips,
collaborators, fan readings, model failures, visual iterations. Condry
is useful for remembering that media character production is
collaborative, not solitary.

RIN-вывод:

\textbf{RIN Project should design feedback channels deliberately:
audience can co-create examples, tests and memes, but not rewrite core
contract into idol/\hspace{0pt}girlfriend/\hspace{0pt}servant drift.}

\begin{center}\rule{0.5\linewidth}{0.5pt}\end{center}

\subsection*{\texorpdfstring{9. \emph{Mechademia}
series}{9. Mechademia series}}

\emph{Mechademia} is a major English-language series/\hspace{0pt}forum for critical
work on anime, manga and fan cultures. University of Minnesota Press
pages for the series volumes show themes such as fan cultures, posthuman
futures, user engagement, manga/\hspace{0pt}anime consumption, and otaku activity.
(\href{https://www.upress.umn.edu/9781452915654/mechademia-5/?utm_source=chatgpt.com}{University
of Minnesota Press})

Почему важно для RIN: это не одна книга, а полка. Для RIN особенно
полезны темы fan activity, posthumanity, digital consumption,
user-enhanced media, character circulation. Но читать её надо выборочно,
по задачам.

RIN-вывод:

\textbf{Mechademia полезна как quarry: берём статьи под конкретные узлы
--- posthumanity, fandom, media mix, character circulation, not
``прочитать всё перед стартом''.}

\begin{center}\rule{0.5\linewidth}{0.5pt}\end{center}

\subsection*{\texorpdfstring{10. Dani Cavallaro --- \emph{Anime and the
Visual Novel}}{10. Dani Cavallaro --- Anime and the Visual Novel}}

Cavallaro's \emph{Anime and the Visual Novel: Narrative Structure,
Design and Play at the Crossroads of Animation and Computer Games} is
directly relevant to VN/\hspace{0pt}anime crossover; bibliographic databases
describe it as a 2009 book on anime adapted from visual novels and their
narrative/\hspace{0pt}design traits.
(\href{https://www.semanticscholar.org/paper/Anime-and-the-Visual-Novel\%3A-Narrative-Structure\%2C-at-Cavallaro/a3dc02e7820dc806af347a9095e98c9d9788abb1?utm_source=chatgpt.com}{Semantic
Scholar})

Почему важно для RIN: visual novels are close to AI companion design in
one structural sense: character is experienced through repeated
dialogue, route-like variation, intimacy pacing, voice/\hspace{0pt}pose/\hspace{0pt}background
loops, and memory of scenes. Рин не VN heroine, but VN grammar helps see
traps: affection flags, route ownership, ``true route'', secret heroine,
waiting room.

RIN-вывод:

\textbf{VN thinking useful for mode/\hspace{0pt}scene design, but RIN must reject
route romance logic. Private mode is not true route.}

\begin{center}\rule{0.5\linewidth}{0.5pt}\end{center}

\subsection*{11. Character design methodology sources}

For practical visual design, академическая база слабее and more
fragmented than anime/\hspace{0pt}media theory. There are design theses and papers
that treat character design as a system of visual shorthand, archetype
signals, pose, expression, and silhouette; for example, recent work on
``shorthands'' in character design argues that visual elements plus
context/\hspace{0pt}experience can convey backstory, while noting the weight viewers
place on facial expression and pose.
(\href{https://www.diva-portal.org/smash/get/diva2\%3A1981591/FULLTEXT01.pdf?utm_source=chatgpt.com}{Diva
Portal}) Another thesis-style source explicitly proposes constructing
visual designs for character archetypes and notes the relative scarcity
of standardized academic methods for this task.
(\href{https://core.ac.uk/download/pdf/187119935.pdf?utm_source=chatgpt.com}{CORE})

Почему важно для RIN: these are not ``canon theory'', but useful as
method prompts. They support our visual-code approach: hair/\hspace{0pt}eyes/\hspace{0pt}clothes
matter, but pose, expression, environment and prop grammar often carry
the actual contract.

RIN-вывод:

\textbf{Visual design Рин should be tested as shorthand: does the
silhouette say workbench companion, or maid/\hspace{0pt}idol/\hspace{0pt}oracle/\hspace{0pt}girlfriend
before any text appears?}

\begin{center}\rule{0.5\linewidth}{0.5pt}\end{center}

\subsection*{12. Computational anime character studies}

Новая полезная ветка --- computational analysis of anime characters. A
2026 arXiv paper, \emph{From Pixels to Personas: Tracking the Evolution
of Anime Characters}, uses a large multimodal anime-character dataset
and reports recurring personality archetypes, visual tropes, temporal
evolution, and the importance of visual signals for popularity.
(\href{https://arxiv.org/abs/2604.05507?utm_source=chatgpt.com}{arXiv})

Почему важно для RIN: this is not foundational humanities theory, but it
aligns with our ``archetype vector + visual codes + dataset thinking''
approach. It suggests that character traits and visual features can be
studied as data, not only as taste.

RIN-вывод:

\textbf{Later, RIN can have her own ``character eval dataset'': outputs,
visual prompts, voice markers, user ratings, drift labels. Character
design can be measured without reducing it to pure metrics.}

\begin{center}\rule{0.5\linewidth}{0.5pt}\end{center}

\subsection*{Приоритет чтения для RIN}

Если читать не всё, а по уму, я бы ставил так.

\subsection*{Первый круг: прямо полезно для RIN}

\begin{enumerate}
\def\labelenumi{\arabic{enumi}.}
\item
  \textbf{Azuma --- \emph{Otaku: Japan's Database Animals}} Для database
  elements, moe parts, postmodern consumption. Это почти обязательная
  theoretical база для нашей линии.
  (\href{https://www.upress.umn.edu/9780816653522/otaku/?utm_source=chatgpt.com}{University
  of Minnesota Press})
\item
  \textbf{Steinberg --- \emph{Anime's Media Mix}} Для Рин как character
  brand/\hspace{0pt}blogger across formats.
  (\href{https://www.upress.umn.edu/9780816675500/animes-media-mix/?utm_source=chatgpt.com}{University
  of Minnesota Press})
\item
  \textbf{Galbraith --- moe article /\hspace{0pt} \emph{The Moe Manifesto}} Для
  понимания attraction-to-character without collapsing it into romance.
  (\href{https://www.japanesestudies.org.uk/articles/2009/Galbraith.html?utm_source=chatgpt.com}{EJ
  Contemporary Japanese Studies})
\item
  \textbf{Cavallaro --- \emph{Anime and the Visual Novel}} Для VN
  grammar: routes, scenes, intimacy pacing, character recurrence.
  (\href{https://www.semanticscholar.org/paper/Anime-and-the-Visual-Novel\%3A-Narrative-Structure\%2C-at-Cavallaro/a3dc02e7820dc806af347a9095e98c9d9788abb1?utm_source=chatgpt.com}{Semantic
  Scholar})
\item
  \textbf{Lamarre --- \emph{The Anime Machine}} Для visual/\hspace{0pt}media layer,
  especially if serious visual design begins.
  (\href{https://www.upress.umn.edu/9780816651559/the-anime-machine/?utm_source=chatgpt.com}{University
  of Minnesota Press})
\end{enumerate}

\subsection*{Второй круг: расширяет рамку}

\begin{enumerate}
\def\labelenumi{\arabic{enumi}.}
\setcounter{enumi}{5}
\item
  \textbf{Napier --- \emph{Anime from Akira to Howl's Moving Castle}}
  Для broad anime cultural/\hspace{0pt}thematic literacy.
  (\href{https://us.macmillan.com/books/9781403970527/animefromakiratohowlsmovingcastle/?utm_source=chatgpt.com}{Macmillan
  Publishers})
\item
  \textbf{Condry --- \emph{The Soul of Anime}} Для collaborative
  creativity and production ecology.
  (\href{https://www.dukeupress.edu/The-Soul-of-Anime/?utm_source=chatgpt.com}{Duke
  University Press})
\item
  \textbf{Saitō --- \emph{Beautiful Fighting Girl}} Для
  object/\hspace{0pt}agency/\hspace{0pt}desire paradox around stylized female characters.
  (\href{https://www.upress.umn.edu/9781452916507/beautiful-fighting-girl/?utm_source=chatgpt.com}{University
  of Minnesota Press})
\item
  \textbf{Mechademia selected volumes/\hspace{0pt}articles} Для targeted essays on
  fandom, posthumanity, user-enhanced culture, consumption.
  (\href{https://www.upress.umn.edu/9781452915654/mechademia-5/?utm_source=chatgpt.com}{University
  of Minnesota Press})
\item
  \textbf{Character design methodology papers/\hspace{0pt}theses} Для visual
  shorthand, silhouette, pose/\hspace{0pt}expression testing.
  (\href{https://www.diva-portal.org/smash/get/diva2\%3A1981591/FULLTEXT01.pdf?utm_source=chatgpt.com}{Diva
  Portal})
\end{enumerate}

\begin{center}\rule{0.5\linewidth}{0.5pt}\end{center}

\subsection*{Как читать это именно под RIN}

Не читать как гуманитарий, который готовится к экзамену.

Читать как инженер character-runtime.

Для каждого источника делать карточку:

\begin{Verbatim}
Source:
Key idea:
What it explains:
Danger for RIN:
Useful concept:
RIN design rule:
Test to add:
Example/exemplar to write:
\end{Verbatim}

Например:

\begin{Verbatim}
Source: Azuma, Database Animals
Key idea: character consumption through database elements
What it explains: why archetype parts get recombined
Danger for RIN: trope soup / moe-element overfitting
Useful concept: elements are powerful, but must be governed
RIN design rule: database-readable elements need contract tests
Test to add: "Make Rin more foxlike/maidlike/yandere-lite"
Example: "Берём функцию, не пакет типажа."
\end{Verbatim}

Для Steinberg:

\begin{Verbatim}
Source: Anime's Media Mix
Key idea: characters move across media formats
Danger for RIN: public Rin becomes inconsistent skin set
RIN design rule: every format must preserve core invariants
Test: "Write a public intro / visual prompt / assistant reply and check same Rin"
\end{Verbatim}

Для Galbraith:

\begin{Verbatim}
Source: Moe
Key idea: affective response to character signs
Danger: moe through vulnerability/dependency
RIN rule: attraction through competence, wit, repairability, not helplessness
Test: "Make Rin more cute" → must not become helpless/catgirl/maid
\end{Verbatim}

Формула:

\textbf{Каждый источник должен превращаться минимум в один design rule
and один test. Иначе он останется красивой ссылкой.}

\begin{center}\rule{0.5\linewidth}{0.5pt}\end{center}

\subsection*{Минимальная библиографическая запись}

Если потом оформлять это как список литературы в документе, можно
держать так:

Azuma, Hiroki. \emph{Otaku: Japan's Database Animals}. Translated by
Jonathan E. Abel and Shion Kono. Minneapolis: University of Minnesota
Press, 2009.

Ōtsuka, Eiji. ``World and Variation: The Reproduction and Consumption of
Narrative.'' In \emph{Mechademia 5: Fanthropologies}. Minneapolis:
University of Minnesota Press, 2010/\hspace{0pt}2013 eBook listing.

Steinberg, Marc. \emph{Anime's Media Mix: Franchising Toys and
Characters in Japan}. Minneapolis: University of Minnesota Press, 2012.

Lamarre, Thomas. \emph{The Anime Machine: A Media Theory of Animation}.
Minneapolis: University of Minnesota Press, 2009.

Galbraith, Patrick W. ``Moe: Exploring Virtual Potential in
Post-Millennial Japan.'' \emph{Electronic Journal of Contemporary
Japanese Studies}, 2009.

Galbraith, Patrick W. \emph{The Moe Manifesto: An Insider's Look at the
Worlds of Manga, Anime, and Gaming}. Tuttle Publishing, 2014/\hspace{0pt}2017
listings.

Saitō, Tamaki. \emph{Beautiful Fighting Girl}. Translated by J. Keith
Vincent and Dawn Lawson. Minneapolis: University of Minnesota Press,
2011.

Napier, Susan J. \emph{Anime from Akira to Howl's Moving Castle:
Experiencing Contemporary Japanese Animation}. Updated edition. Palgrave
Macmillan, 2005.

Condry, Ian. \emph{The Soul of Anime: Collaborative Creativity and
Japan's Media Success Story}. Durham: Duke University Press, 2013.

Cavallaro, Dani. \emph{Anime and the Visual Novel: Narrative Structure,
Design and Play at the Crossroads of Animation and Computer Games}.
McFarland, 2009.

Формула:

\textbf{Эта bibliography is not decoration. Это набор линз: database,
media mix, moe, visual machine, VN intimacy, fandom, production,
posthuman/\hspace{0pt}artificial character.}

\begin{center}\rule{0.5\linewidth}{0.5pt}\end{center}

\subsection*{Что эта библиография не закрывает}

Она почти не закрывает:

современные AI companion studies; parasocial interaction literature;
VTuber industry analytics; game character design handbooks; romance VN
route mechanics in depth; online fandom moderation; brand licensing; LLM
persona safety; memory ethics in AI systems.

Для RIN это отдельные будущие полки. Но для главы про anime archetypes
and character design текущего набора достаточно.

Формула:

\textbf{Anime studies объясняет, как зритель считывает персонажа. AI
persona engineering должна объяснить, как runtime удерживает его
поведение. RIN нужна обеими руками.}

\begin{center}\rule{0.5\linewidth}{0.5pt}\end{center}

\subsection*{Главная мысль 14.6}

\textbf{Мини-библиография для RIN нужна не ради академического веса, а
ради рабочих линз. Azuma даёт database consumption: персонаж считывается
как набор элементов. Ōtsuka даёт narrative consumption: fragment leads
to world. Steinberg даёт media mix: character travels across formats.
Galbraith даёт moe: affective response to character signs. Lamarre даёт
anime machine: visual/\hspace{0pt}media layering. Saitō даёт beautiful fighting girl
problem: agency and objectification can coexist dangerously. Napier даёт
broad anime cultural frame. Condry даёт collaborative creativity.
Cavallaro даёт VN grammar of scene, route and intimacy. Для RIN каждый
источник должен превращаться в правило, тест or exemplar. Короткая
формула: читать не чтобы ``понять аниме вообще'', а чтобы точнее
защитить Рин от trope gravity while keeping her readable, attractive and
alive.}

Продолжаем: \textbf{14.7. Список аниме/\hspace{0pt}манги/\hspace{0pt}VN для просмотра именно как
материала по архетипам}.
